\section{Alternative explanations}
\label{sec:alternative_explanations}

In this section, we conduct several robustness exercises. We first show the robustness of our results to a rich set of controls. We then show that financial sector distress did not materially impact public firms and provide further evidence that financial frictions do not explain our results. We also show that the industry-specific policies of the National Industrial Recovery Act (NIRA) cannot explain the investment gap we document. We then establish that the dollar's devaluation did not differentially impact firms based on their exposure to gold-denominated debt. Lastly, we show that our results are robust to a wide range of alternative specifications.

\subsection{Controlling for firm characteristics}
\label{sec:control}

To address the concern that pretreatment differences in firm characteristics between the control and treatment groups reflect differential exposure to macreoeconomic and investment opportunity shocks, we run the following regression:
\begin{equation}
\label{eq:control}
\text{Net investment}_{i,t} = \beta_0 + \beta_1 d_i p_t +  \sum_j  \sum_{\tau} \beta_{j,\tau}  \mathbb{I}_{t = \tau} X_{j,i} + \alpha_i  + \delta_t+ u_{i,t}, 
\end{equation}
where $X_{j,i}$ is firm $i$'s pretreatment value for characteristic $j$. In each regression, $X_{j,i}$ is computed as the average value of the firm-level variable in the pretreatment period with $p_t = 0$. The pretreatment characteristics $X_{j,i}$ are interacted with time indicator functions for each year, denoted by $\mathbb{I}_{t = \tau}$, in order to account for variation in the impact of observable characteristics on investment over time. Column (1) of Table \ref{tab:T5} reports results where the set of $X_{j,i}$ includes log(Assets), Tobin's Q, book leverage, financial leverage, cash/assets, net income/assets, and noncurrent assets/assets, which we found to be significantly different between $d = 0$ and $d > 0$ firms on average. We find that a statistically significant estimate of $\beta_1$ is -4.6\% (Panel A) for the emergence of leverage risk and 6.3\% (Panel B) for the elimination of leverage risk. This is economically and statistically similar to our findings  confirming our conjecture that pretreatment differences do not explain the impact of leverage risk on investment.  

We also verify that our results are robust to the inclusion of nonlinear controls, and we reestimate equation (\ref{eq:control}) using decile dummies for each characteristic.\footnote{In this case, $X_{j,i}$ is an indicator function whether firm $i$ belongs to a particular decile of the pretreatment distribution of the characteristics and $j \in \{1,..., 10\}$.} Doing so greatly increases the number of regressors, and, because of the size of our sample, we include in the regression one control variable at a time. Columns (2) to (8) of Table \ref{tab:T5} report the results. Overall, including a nonlinear control does not significantly change our coefficient estimates. Hence, despite differences in the average pretreatment values of firm attributes, there is variation in $d_i$ affecting the investment path among otherwise similar firms.

\subsection{Credit constraints}
\subsubsection{Bank loans}
In his seminal paper, \cite{Bernanke1983} argues that the disruption in the banking sector during the Great Depression reduced the supply of credit available to borrowers. Although the contraction in bank credit might explain why smaller firms reduced their investment, this channel is unable to explain the investment gap related to firms' gold clause exposure for at least two reasons.

First, as \cite{Bernanke1983} acknowledges, ``most larger corporations entered the decade with sufficient cash and liquid reserves to finance operations and any desired expansions.'' Consistent with this assessment, public firms with gold-denominated debt in our sample had about 10\% of their assets in cash and liquid assets throughout the 1931 to 1936 period. This amount of cash would have been sufficient to finance the investment gap caused by leverage risk for about 3 years for firms with $d = 1$. Second, public firms rarely relied on bank debt as a source of financing. As shown in Table \ref{tab:T1}, the sum of bank loans and notes payable on average account for 3\% of the total liabilities of these firms. This sum is an order of magnitude smaller than the contribution of corporate bonds (44\%). 

Nonetheless, we test whether firms' reliance on bank debt can explain the investment gap by replacing our leverage risk exposure measure $d_i$ by the share of bank debt in total liabilities. Specification (1) of Table \ref{tab:T7} shows the lack of evidence that this alternative measure of exposure explains cross-sectional variation in firms' investment.

\subsubsection{Corporate bonds}

\cite{Benmelech2018} show that firms with corporate bonds maturing during the Great Depression until 1934 were subject to more severe financial frictions, since the collapse of the corporate bond primary market during the crisis made it difficult to rollover maturing debt. Because firms with more gold-denominated corporate bonds are more likely to have a bond issue maturing during the crisis (or at any point in time), a concern is that our leverage risk measure $d_i$ is a proxy for the credit constraints of firms with maturing bond issues. 

To rule out this possibility, we collect data on the maturity date of the bond issues and remove from the sample any firm with an outstanding bond maturing in 1933 or 1934. The procedure eliminates 29 firms from the 1931--1934 sample and 30 firms from the 1933--1936 sample. Specification (2) of Table \ref{tab:T7} reports that our main results remain unaffected using this alternative sample. Leverage risk matters even for firms that did not have bonds maturing during the Great Depression.

\subsection{New Deal policies}

Following the inauguration of FDR in early 1933, a series of programs, public projects, financial reforms, and regulations were enacted. Collectively referred to as the New Deal, the most important initiatives included the Emergency Banking Act, the Glass-Steagall Act, the Agricultural Adjustment Act (AAA), and the National Industrial Recovery Act (NIRA). The effects of these reforms on the economy are still debated to this day, with some economists suggesting that the New Deal policies kept the economy depressed beyond 1933 (see \cite{Cole2004} for a discussion of the New Deal and its impacts). If firms with more gold-denominated debt were incidentally more affected by New Deal policies, the association between leverage risk and investment might be related to these policies.

For instance, \cite{Cole2004} argue that the industrial codes included in the NIRA created labor market rigidities that kept the economy depressed. The NIRA regulations prescribed specific operating rules that applied to all firms within an industry, but some industries faced tighter restrictions than did others. Similarly, the AAA introduced regulations specific to the agricultural sector, and the two banking acts were directed at the financial sector.\footnote{Financial firms are excluded from our analysis, and, as previously discussed, the public firms included in our sample did not significantly rely on bank debt as a source of financing.} To control for the potentially confounding effects of the New Deal policies on our results, we add 2-digit SIC industry-year fixed effects to equation (\ref{eq:main}) and reestimate the model.

Table \ref{tab:T8} reports the results. Including industry-year fixed effects to the regression model does not affect our results significantly. Specification (1) of Panel A shows that, when industry-year fixed effects are included, the estimated treatment effect coefficient is -5.4\% (compared to -5.9\% without). Similarly, specification (1) of Panel B shows that the estimated treatment reversal coefficient is 4.9\% when industry-year fixed effects are included (compared to 5.1\% without). Column (2) of Table \ref{tab:T8} confirms that the results with industry-year fixed effects are robust to the inclusion of firm-level controls in the regression. Overall, the results suggest that leverage risk remains significant even after controlling for the industry-wide effects of the New Deal policies.

\subsection{Dollar's devaluation and the farm channel}

In their study of the early stages of the recovery, \cite{Hausman2017} show that the dollar's devaluation in 1933 led to a large increase in traded crop prices and to a higher demand for goods in the agricultural regions of the country because of the increased purchasing power of farmers. If firms with little exposure to gold-denominated debt incidentally benefited more from the rise in demand, they would have higher sales and profits explaining why those firms invested more.

To determine whether the investment gap we document could be explained by demand shocks, we reestimate equation (\ref{eq:main}) with profits normalized by total assets as the dependent variable.\footnote{We use profits rather than sales, because most income statements in Moody's Manuals directly report net income rather than sales and costs separately.} Table \ref{tab:T9} reports the results. There is no significant relationship between the gold clause content of a firm's liabilities and profits. It is thus unlikely that our results can be attributed to a positive demand shock specific to firms without exposure to gold clauses. 

\subsection{Other robustness checks}

We experiment with various alternative measures of investment and exposure to gold clauses to show the robustness of our results in Columns (1) to (5) of Table \ref{tab:TA1}. In Column (1), we define book debt net of preferred shares. In Column (2), we use the log investment rate as the dependent variable. In Column (3), we omit other regulated industries, namely, transportation, communication, and utilities. In Column (4), we define book debt net of cash. In Column (5), we define $d_i$ using total assets as the denominator instead of total liabilities. Both our results for the emergence (Panel A) and elimination (Panel B) of leverage risk are robust to all alternative specifications.

To ensure that our results are not driven by firms' differential exposure to macroeconomic shocks related to the duration of firms' liabilities, we compute the share of preferred shares and bonds with no gold clauses in total liabilities.\footnote{See \cite{Diamond2014} for a theoretical discussion of the relation between liability duration and debt overhang.} This ratio represents the share of long-term liabilities that are not exposed to leverage risk. Specification (6) in Table \ref{tab:TA1} shows that the share of other long-term liabilities does not generate any dispersion in investment rates upon the emergence of leverage risk in 1933 or its elimination in 1935.